% Generated draft: requires mathematical and source review.
\documentclass[11pt]{article}
\usepackage[margin=25mm]{geometry}
\usepackage{fontspec}
\setmainfont{Latin Modern Roman}
\usepackage{amsmath,amssymb,mathtools,unicode-math}
\setmathfont{Latin Modern Math}
\usepackage{enumitem,needspace,xurl,hyperref}
\setlength{\parindent}{0pt}
\setlength{\parskip}{5pt}
\setlength{\emergencystretch}{3em}
\newcommand{\N}{\mathbb{N}}
\newcommand{\Z}{\mathbb{Z}}
\newcommand{\Q}{\mathbb{Q}}
\newcommand{\R}{\mathbb{R}}
\newcommand{\C}{\mathbb{C}}
\newcommand{\F}{\mathbb{F}}
\newcommand{\A}{\mathbb{A}}
\newcommand{\PP}{\mathbb{P}}
\newcommand{\E}{\mathbb{E}}
\newcommand{\eps}{\varepsilon}
\newcommand{\dd}{\,\mathrm{d}}
\begin{document}
\section*{UnsolvedMath ID: 20002370}
\textbf{Source problem code:} catalogue entry
\subsection{Definitions and mathematical statement}
Correct the OCR notation by putting \(\R_{\mathrm{alg}}=\overline{\Q}\cap\R\), the field of real algebraic numbers. Let
\[
 K=\R_{\mathrm{alg}}(t),\qquad L=\R(t),
\]
where \(t\) is transcendental. The intended language in the supplied aid is the parameter-free pure ring language \(\{0,1,+,-,\cdot\}\); in particular there is no named constant for \(t\). The source states that \(\R_{\mathrm{alg}}\equiv\R\), whereas \(K\not\equiv L\), and that \(K\) and \(L\) have the same existential theory. Zahidi's Question 34 asks for the least quantifier complexity of a sentence distinguishing \(K\) and \(L\), and separately the least number of variables needed.

For clarity, distinct complexity measures must not be conflated. For a prenex sentence, its block-alternation depth is the number of changes between existential and universal blocks, so a single quantifier block has depth zero. A \(\forall_n\exists\) sentence means
\[
 \forall x_1\cdots\forall x_n\ \exists y_1\cdots\exists y_m\,
 \theta(x_1,\ldots,x_n,y_1,\ldots,y_m),
\]
for some finite \(m\), with \(\theta\) quantifier free. The leading-block size \(n\), total quantifier count, quantifier rank, and number of distinct object-variable symbols in an arbitrary formula are different measures.

The supplied partial-result report states that \(K\) and \(L\) agree on all existential sentences, all universal sentences, all \(\forall_1\exists\) sentences, and their dual \(\exists_1\forall\) sentences, allowing arbitrary finite trailing blocks and quantifier-free Boolean matrices. It also reports a published Liouville-cut construction yielding a parameter-free pure-ring separator of the form
\[
 \forall x_1\forall x_2\forall x_3\ \exists y_1\cdots\exists y_m\,
 \theta(x_1,x_2,x_3,y_1,\ldots,y_m)
\]
for some finite \(m\). Neither the matrix, its orientation between \(K\) and \(L\), nor the number \(m\) is supplied.

Subject to these reported results, the minimum prenex block-alternation depth among separators is exactly one. For the more specific quantity
\[
 n_{\min}=\min\{n:\text{some parameter-free }\forall_n\exists\text{ sentence distinguishes }K,L\},
\]
the aid gives only
\[
 2\le n_{\min}\le3.
\]
It does not decide whether two leading universal variables suffice. Finally, if \(v_{\min}\) is the least number of distinct object-variable symbols in any parameter-free ring sentence distinguishing the fields, with reuse and rebinding of symbols allowed, the reported lower bound is \(v_{\min}\ge2\); its exact value remains unresolved in the aid. The three-variable leading universal block does not bound the total number of variables by three, because the existential witnesses also require variables. No exact minimum total quantifier count or quantifier rank is asserted.

All these bounds are literature and partial-result assertions supplied by the record, not a separator constructed or verified here. Their language restrictions are essential; changing the language by adding parameters or constants is a different complexity question.
\subsection{Short English statement}
For the two real rational-function fields in the pure ring language, the aid reports an exact minimum of one quantifier alternation. It leaves the least leading universal block at two or three, and the least total number of variable symbols unresolved beyond a lower bound of two.
\subsection{Sources}
\begin{itemize}
\item Local registry, permanent UnsolvedMath ID \(20002370\), exact public code \texttt{catalogue entry}; canonical location \texttt{aim-logic-notes.json}, \texttt{notes[145]}.
\item AIM, \emph{Extensions of Hilbert's tenth problem}, Question 34 attributed to Zahidi, page 7 according to the supplied verification field; reference \url{https://aimath.org/WWN/hilberts10th/hilberts10th.pdf}. The verification field reports that OCR lost the bar in \(\overline{\Q}\) and rendered \(\not\equiv\) incorrectly. No source inspection was performed here. No proposer year, original remarks, or literature entry is supplied.
\item Supplied research attempt 1: \texttt{partial\_result}, run and registry status \texttt{partially\_solved}, rigor label \texttt{proved}, low novelty confidence; classification \texttt{PARTIAL-PROGRESS}, created and updated 14 August 2026. It reports one-universal-variable preservation and a published Liouville-cut separator, but supplies no author, title, journal, theorem number, or explicit separating sentence.
\item The research aid is explicitly machine generated, not peer reviewed, and requires independent verification before citation.
\end{itemize}
\end{document}
